\section{Policy specifications}
\label{app:policies}

This appendix pins down the formal pieces that Section~\ref{sec:agents} summarises: the typed tool surface, the statistical primitives, the system prompts, the $z$-calibrated mean-shift rule, the per-policy allowed-action sets, and the runtime loop. The policy comparison table itself is not restated. Code paths refer to \texttt{run\_ladder.py} and the modules under \texttt{src/causal\_discovery/agents/}.

\subsection{Tool surface}
\label{app:policies:tools}

Each turn emits exactly one action via a strict JSON-schema tool call: the schema is generated by \texttt{make\_action\_response\_schema} (\texttt{tool\_schema.py}) with \texttt{strict: True} and \texttt{additionalProperties: False}, so every required field for the chosen action must be present, unused scalar fields are \texttt{null}, and unused list fields are \texttt{[]}. The five action signatures, lifted verbatim from \texttt{ACTION\_FIELDS} and \texttt{FIELD\_SCHEMAS}, are:

\begin{center}
\small
\begin{tabular}{@{}l p{0.72\linewidth}@{}}
\toprule
Action & Signature \\
\midrule
\texttt{intervene}            & \texttt{(var: int, value: float)} \\
\texttt{correlation}          & \texttt{(i: int, j: int) -> float} \\
\texttt{partial\_correlation} & \texttt{(i: int, j: int, conditioning\_on: list[int]) -> float} \\
\texttt{independence\_test}   & \texttt{(i: int, j: int, conditioning\_on: list[int], alpha: float) -> (bool, float)} \\
\texttt{submit\_graph}        & \texttt{(directed\_edges: list[[int,int]], undirected\_edges: list[[int,int]], reasoning\_summary: str)} \\
\bottomrule
\end{tabular}
\end{center}

Every action also carries a \texttt{reasoning\_summary} string, which is appended to a working-memory deque (Appendix~\ref{app:policies:loop}).

\subsection{Statistical primitives}
\label{app:policies:stats}

The three CI-style tools are implemented in \texttt{stats\_tools.py} and operate on the \emph{active dataset}: the observational sample $D^{\text{obs}}$ before any intervention, and the most recent intervention sample thereafter (\texttt{run\_ladder.py:578--579, 602--603, 629--630}).

\begin{itemize}
\itemsep0pt
\item \texttt{correlation(i, j)} returns the Pearson correlation $\hat\rho_{ij}$ of columns $i$ and $j$ of the active dataset, clipped to $[-1, 1]$.
\item \texttt{partial\_correlation(i, j, Z)} regresses $X_i$ and $X_j$ on $[\mathbf{1}, X_Z]$ via OLS (\texttt{numpy.linalg.lstsq}) and returns the Pearson correlation of the residuals; with $Z = \emptyset$ it reduces to \texttt{correlation(i, j)}.
\item \texttt{independence\_test(i, j, Z, $\alpha$)} performs the Fisher-$z$ test:
\begin{equation}
\hat\rho = \mathrm{partial\_correlation}(i, j, Z), \qquad
z = \tfrac{1}{2}\log\frac{1 + \hat\rho}{1 - \hat\rho}\sqrt{n - |Z| - 3}, \qquad
p = \operatorname{erfc}(|z|/\sqrt{2}),
\end{equation}
where $\hat\rho$ is first clipped to $\pm(1 - 10^{-6})$ for numerical stability. The tool returns the pair \texttt{(p > $\alpha$,\, p)}, declaring conditional independence when $p > \alpha$. The test errors out when $\mathrm{dof} = n - |Z| - 3 \le 0$.
\end{itemize}

\subsection{System prompts}
\label{app:policies:prompts}

Two system prompts are used. The \emph{raw} prompt (\texttt{build\_system\_prompt\_raw}) tells the model that it is performing causal discovery over an unknown linear-Gaussian DAG with full observability and no hidden confounders; cautions that dependence and correlation are not direct causation by themselves, so chains and shared causes can induce spurious dependence; defines an intervention as forcing one variable to a fixed value, breaking its incoming dependencies while leaving other relationships intact; suggests choosing intervention values far enough from the observational mean to make downstream shifts detectable while staying on the observed scale; lists the allowed actions; and instructs the model to use the action tool exactly once per turn and not to request hidden metadata.

The \emph{stats} prompt (\texttt{build\_system\_prompt\_stats}) uses the same base and adds three blocks. A \emph{tool purpose} block: \texttt{correlation} is a coarse association check that is not causal direction by itself, \texttt{partial\_correlation} tests direct dependence after conditioning, and \texttt{independence\_test} decides whether evidence supports conditional independence. A \emph{tool data source} block: tools use the observational dataset before any intervention and the latest intervention dataset after one, and may return \texttt{value=null} with an error when a statistic is undefined. A \emph{decision policy} block: gather evidence, orient edges when the evidence supports direction, and leave undirected only when evidence is ambiguous; in observational-only configurations the model is also told to perform at least three statistical actions before its first \texttt{submit\_graph}.

The session prompt (\texttt{build\_session\_prompt}) is a JSON object with keys \texttt{variables}, \texttt{budget}, \texttt{output\_schema}, \texttt{json\_contract}, optionally \texttt{observational\_data}, and any additional \texttt{session\_context} entries; the contract requires every output-schema field to appear on every turn, with \texttt{null} for unused scalar fields, \texttt{[]} for unused list fields, and a brief \texttt{reasoning\_summary}.

\subsection{Observation format}
\label{app:policies:observation}

The model's input is \emph{numerical}: the observational sample is delivered as a $n_\text{obs} \times d$ JSON array of raw floating-point values, not as natural-language descriptions, summary statistics, or QA pairs. Variables are referred to by index-encoded names \texttt{X0, X1, \ldots, X(d-1)} carrying no domain semantics; there is no metadata that hints at the underlying graph. Every turn rebuilds the session-prompt JSON object from scratch with the following keys:

\begin{itemize}
\itemsep0pt
\item \texttt{variables}: list of $d$ string names \texttt{["X0", "X1", \ldots]}.
\item \texttt{budget}: integer count of remaining intervention turns (decremented after each \texttt{intervene}).
\item \texttt{output\_schema} / \texttt{json\_contract}: per-action field requirements as in §\ref{app:policies:prompts}.
\item \texttt{observational\_data} (when \texttt{include\_observational\_data=True}): $n_\text{obs} \times d$ list of lists of float, full machine precision. Withheld for \texttt{pc\_cpdag\_llm}, where the model receives only the PC partial graph and an \texttt{observational\_summary} of per-variable means and standard deviations to three decimals (§\ref{app:policies:perpolicy}).
\item \texttt{intervention\_history} (after the first \texttt{intervene}): list of \texttt{\{step, var, value, data, remaining\_budget\}} entries, where \texttt{data} is the $n_\text{int} \times d$ raw float matrix from that intervention.
\item \texttt{working\_memory}: the deque of prior \texttt{reasoning\_summary} strings (max $12$, $\le 300$ chars each; §\ref{app:policies:loop}).
\end{itemize}

For the canonical L3 instance of Appendix~\ref{app:traces:paired} ($d = 10$, $n_\text{obs} = 50$, $n_\text{int} = 25$), the step-1 user message is roughly $5$~kB of JSON. After one intervention the step-2 message expands to roughly $15$~kB as the $25 \times 10$ intervention sample is appended to \texttt{intervention\_history}; subsequent interventions concatenate further blocks. There is no token-budget pruning: the model sees the entire numerical history every turn. The first $\le 250$ characters of the step-1 \texttt{observational\_data} value are reproduced verbatim below to make the format unambiguous (full $50 \times 10$ payload available in the on-disk \texttt{events.jsonl} per Appendix~\ref{app:repro:repo}):

\begin{center}
\small
\begin{verbatim}
"observational_data":[[-0.5996848527,0.1257563716,0.3868730550, ... ,1.1959564656],
                      [ 0.2278882859,0.2679279303,2.1911019666, ... ,-0.2995619661],
                      [-0.2708565131,1.9276286132,-0.7202341092, ... , 0.1743317379],
                      ...  // 47 more rows
                     ]
\end{verbatim}
\end{center}

This design makes the benchmark a test of \emph{structure recovery from raw numerical evidence}: the model must perform the regression / mean-shift / partial-correlation reasoning itself (or invoke the statistical primitives of §\ref{app:policies:stats} when allowed), not pattern-match against a pre-summarised description. Appendix~\ref{app:traces} shows full agent traces operating on this format.

The raw-numerical format is a deliberate design choice, not a convenience. Pre-summarised inputs --- correlation tables, natural-language descriptions of dependencies, or rounded statistics --- would conflate causal reasoning with reading comprehension, and would make the \texttt{llm\_raw}, \texttt{llm\_stats}, and \texttt{pc\_cpdag\_llm} contrasts of §\ref{app:policies:perpolicy} meaningless: every policy would already start from a summary, so the score gap between them would no longer measure what the model extracts from raw evidence. By delivering raw samples, the benchmark forces the model to perform the inference itself when allowed only \texttt{intervene}/\texttt{submit\_graph}, and to choose \emph{which} statistical queries to issue when the stats tools are available; the resulting policy gaps reported in Appendix~\ref{app:results} measure exactly that increment of value. The cost of including the full numerical history every turn (§\ref{app:results:cost}) is reported alongside the score so a reader can judge the dollar-per-F1-point tradeoff this design imposes.

\subsection{\texorpdfstring{$z$-calibrated mean-shift rule}{z-calibrated mean-shift rule}}
\label{app:policies:zrule}

The classical resolver (\texttt{resolve\_with\_interventions}, \texttt{run\_ladder.py:278--314}) takes a partially directed graph with directed edge set $E_d$ and undirected edge set $E_u$, the observational sample $D^{\text{obs}}$, and an environment with a remaining intervention budget. It iterates the following three steps while $E_u$ is non-empty and the budget is positive.

\paragraph{Target selection.} Let $\deg_U(j)$ be the degree of node $j$ in the current undirected subgraph. Pick
\begin{equation}
j^* = \arg\min_{j \in V_U}\, \bigl(-\deg_U(j),\; j\bigr),
\end{equation}
i.e.\ the node of highest undirected degree, with ties broken by lowest node index (\texttt{run\_ladder.py:289--293}).

\paragraph{Intervention.} Set $v = \hat\mu_{\text{obs}}[j^*] + \delta$ with $\delta = 3.0$ (\texttt{INTERVENTION\_OFFSET}, \texttt{run\_ladder.py:56}) and call \texttt{env.intervene}$(j^*, v)$ to draw $n_{\text{int}}$ samples.

\paragraph{Orientation.} For each undirected edge $\{j^*, k\} \in E_u$, compute the threshold
\begin{equation}
\tau_k = z_{0.975}\, \sqrt{\frac{\hat\sigma^2_{\text{obs}}[k]}{n_{\text{obs}}} + \frac{\hat\sigma^2_{\text{int}}[k]}{n_{\text{int}}}}, \qquad z_{0.975} = 1.959963984540054,
\end{equation}
where the variances are sample variances with $\mathrm{ddof} = 1$ (\texttt{run\_ladder.py:285, 298}) and $z_{0.975}$ is \texttt{MEAN\_SHIFT\_Z\_975} (\texttt{run\_ladder.py:55}). Orient $j^* \to k$ if $\lvert\hat\mu_{\text{int}}[k] - \hat\mu_{\text{obs}}[k]\rvert > \tau_k$, otherwise $k \to j^*$. If the chosen orientation would create a directed cycle in $E_d$, the edge is left undirected and the loop continues (\texttt{run\_ladder.py:310--311}, using \texttt{would\_create\_cycle}).

\subsection{Per-policy specification}
\label{app:policies:perpolicy}

The four LLM policies share the loop in Appendix~\ref{app:policies:loop}; they differ only in the allowed-action set (\texttt{allowed\_actions\_for\_method}, \texttt{tool\_schema.py:102--128}), the system prompt, the session-prompt context, and the action budget (\texttt{--max-steps-raw}, default $20$; \texttt{--max-steps-stats}, default $40$; \texttt{run\_ladder.py:992--993}).

\paragraph{\texttt{pc\_greedy} (\texttt{run\_ladder.py:436--453}).} No LLM. The PC algorithm with Fisher-$z$ conditional-independence tests at $\alpha = 0.05$ runs on $D^{\text{obs}}$ via \texttt{causallearn}, producing an initial directed/undirected partition; the resolver of Appendix~\ref{app:policies:zrule} then consumes the full intervention budget.

\paragraph{\texttt{llm\_raw}.} Allowed actions \{\texttt{intervene}, \texttt{submit\_graph}\}; raw prompt; the raw observational matrix is included in the session prompt; budget \texttt{--max-steps-raw}.

\paragraph{\texttt{llm\_stats}.} Allowed actions \{\texttt{correlation}, \texttt{partial\_correlation}, \texttt{independence\_test}, \texttt{intervene}, \texttt{submit\_graph}\}; stats prompt; same observational data as \texttt{llm\_raw}; budget \texttt{--max-steps-stats}.

\paragraph{\texttt{pc\_cpdag\_llm} (\texttt{run\_ladder.py:727--772}).} PC runs first on $D^{\text{obs}}$. The LLM session is then initialised with \texttt{include\_observational\_data=False} --- the raw matrix is withheld --- and a \texttt{session\_context} containing the PC \texttt{initial\_directed\_edges} and \texttt{initial\_undirected\_edges}, an \texttt{observational\_summary} with per-variable means and standard deviations rounded to three decimals, and an instruction to start from the supplied PC partial graph and use interventions to resolve uncertain directions. Allowed actions \{\texttt{intervene}, \texttt{submit\_graph}\}; raw prompt; budget \texttt{--max-steps-raw}.

\paragraph{\texttt{llm\_stats\_cpdag\_greedy} (\texttt{run\_ladder.py:775--811}).} Allowed actions \{\texttt{correlation}, \texttt{partial\_correlation}, \texttt{independence\_test}, \texttt{submit\_graph}\} --- \texttt{intervene} is removed; stats prompt; budget \texttt{--max-steps-stats}. After the LLM submits, the resolver of Appendix~\ref{app:policies:zrule} runs post-hoc on the LLM's $E_d$ and $E_u$ against the remaining intervention budget; the additional intervention count is added to \texttt{intervene\_actions} in the run diagnostics.

\subsection{Loop semantics}
\label{app:policies:loop}

LLM episodes execute as a stateful sequential loop. The session prompt is rebuilt every turn from the current \texttt{remaining\_budget} and the (possibly coerced) effective allowed-action set; tool history accumulates across turns; the working-memory deque keeps at most $12$ reasoning summaries of $\le 300$ characters each (\texttt{llm.py:28--29}). On the final step the runtime forces \texttt{effective\_allowed\_actions = \{submit\_graph\}} regardless of the policy's nominal action set, so the episode always terminates with a submission (\texttt{run\_ladder.py:531--532}). At submission time, undirected edges that overlap a submitted directed edge are silently dropped before scoring; the dropped count is logged as \texttt{sanitized\_overlap\_count} (\texttt{run\_ladder.py:259--275, 664--668}). This enforces the disjointness invariant of Appendix~\ref{app:scoring} even when the model violates it.
